\chapter{Frontiers}

\section{What this chapter is}

The book has been building a vocabulary. By Chapter~8 the vocabulary
was complete: a five-slot lens for reading any artificial-life system,
a worked composite example showing how the slots compose, a
measurement framework for emergence and novelty. This chapter is about
the questions the vocabulary does not yet answer.

Some of these are technical conjectures the framework has stated and
deferred. Some are open problems that would require new mathematics.
Some are domains where the lens has not yet been seriously applied.
And some are questions that are not really technical at all: questions
about what artificial-life research is for, and where it should go.

The chapter is organised around the conjectures and open problems
already named in the paper (and tracked in \texttt{RESULTS.md}). After
those, a few looser thoughts about where the lens might travel.

\section{C1: closure-operator unification of viability}

We met this in Chapter~6. The four viability formalisms ---
Markov-blanket, autopoietic closure, RAF set, and Minimal Criterion
Coevolution --- all share the rough shape of ``a supporting set is
viable iff it is closed under some relation.'' C1 conjectures that
each instantiates a closure operator on a poset of supporting sets,
and that the four formalisms are the same notion seen four ways.

The status: two cases (autopoietic closure, RAF) are essentially
proved. The Markov-blanket case is genuinely open --- the statistical
decomposition into internal/blanket/external states is not obviously
a closure operator. The MCC case is the hardest --- MCC is a
predicate, not a closure, and the reformulation work has not been
done.

A proof of C1 would unify the four formalisms abstractly. A
counterexample --- say, a system that passes a Markov-blanket
viability test but whose supporting set is provably not closed under
any natural closure on its substrate --- would reveal that
``viability'' is genuinely heterogeneous in the same way ``life'' is
heterogeneous across biological taxa.

The book deliberately ships the four formalisms as separate
protocols in the Python scaffold. This is not laziness; it is
sceptical optimism. If C1 is later proved, the scaffold can be
refactored without changing user code. If it is refuted, the
scaffold has lost nothing.

\section{C2: hierarchical viability composition}

Viability formalisms apply at one level. Real systems are nested:
cells in organisms, ants in colonies, individuals in cultures. C2
conjectures that each viability formalism lifts to hierarchical
entities such that $F^*$-validity at the composite level implies
$F$-validity of each component.

The status: Beer's autopoietic gliders compose correctly under
simple unions. RAFs do not in general --- the union of two RAFs is
not necessarily a RAF without their joint catalyses. This is a
partial counterexample for one direction of C2.

A productive next step is empirical: identify a Lenia-based system
where hierarchical autopoietic closure is non-trivial (colonies of
creatures whose collective closure is more than the closure of
individuals), and check whether the composite-level closure conditions
hold by simulation. The technical apparatus exists; the
implementation does not.

\section{C3: observer-relative novelty as divergence}

Chapter~7 introduced three observer families: density-distance,
foundation-model embedding, and information-theoretic. C3 conjectures
that all three are instances of a single divergence-functional schema
$O[\mu] = D(\mu \,\|\, \mu_{\mathrm{ref}})$ with context-dependent
reference $\mu_{\mathrm{ref}}$.

The status: plausible by inspection, unproven in generality.
Density-distance fits by construction. FM-embedding fits if you let
$\mu_{\mathrm{ref}}$ be the archive's embedding distribution.
Information-theoretic fits if you let $\mu_{\mathrm{ref}}$ be the
appropriate factored distribution.

C3 is the most likely to be true and the least consequential if it
is. Even if proved, the three families remain distinct in practice
(different reference distributions, different computational profiles,
different reproducibility requirements). The proof would clean up the
classification but not change much else.

\section{OP1: causal emergence of scaffolded systems}

This is the framework's biggest open problem. The Hoel definition of
causal emergence applies to any Markov chain; the framework's $M(
\mathcal{S})$ is a particular kind of Markov chain. OP1 asks: for
which choices of substrate, variation, viability, topology, and
observer does $\sup_\Pi \mathrm{CE}(M(\mathcal{S}), \Pi) > 0$?

The question has three sub-questions:

\begin{itemize}[leftmargin=1.5em,itemsep=2pt]
  \item \emph{Are Flow-Lenia substrates with Markov-blanket viability
    causally emergent?} Plausibly yes, but unproved.
  \item \emph{Are differentiable-logic CAs implementing Conway's Game
    of Life causally emergent?} Israeli and Goldenfeld showed Class
    III/IV ECAs admit non-trivial coarse-grainings; does the GoL
    follow?
  \item \emph{Is POET / Enhanced POET causally emergent at the
    environment-agent level?} The dynamics is complex; is there a
    coarse-graining at which environment-agent dynamics achieve
    higher EI than the joint microscale?
\end{itemize}

OP1 is concrete and tractable in pieces. Pick a small candidate
system, build the Markov chain, search for high-EI coarse-grainings.
Whether the framework's apparatus genuinely produces causal emergence
for any of its representative cases is the most important empirical
question the book can pose.

\section{OP2: FEP-as-lens scope}

The free energy principle is treated in the framework as a
descriptive lens (Chapter~2) rather than a derivation. OP2 asks:
characterise precisely the substrate--dynamics pairs for which the
FEP is informative as a parameterisation of viability and observer
slots.

This is partly a technical question (when does sparse coupling
imply genuine Markov-blanket structure?) and partly a methodological
question (what is the FEP \emph{adding} when it applies?). Both halves
are open.

A useful sub-problem: for the Flow-Lenia + LLM system of Chapter~8,
is the Markov-blanket viability filter $F_{\mathrm{MB}}$ a
\emph{strictly stronger} predicate than the autopoietic closure
filter $F_{\mathrm{AC}}$? If yes, $F_{\mathrm{MB}}$ excludes some
viable systems; the FEP is informative about which ones. If no, the
two formalisms are interchangeable for this case and the FEP is
descriptively redundant. The book conjectures the answer depends on
the system; OP2 asks for the characterisation.

\section{OP3: interaction-topology expressiveness}

Chapter~7 introduced the multiplex with $\diamond_\rho$ coupling as
the working answer to multi-regime interactions. OP3 asks whether
this formalism is universal: can every ``natural'' multi-coupling
ALife system be expressed as a multiplex with countably many layers?

A categorical formulation would express OP3 in terms of the operad of
wiring diagrams. The conjecture, in that language: the category of
multiplex topologies with $\diamond_\rho$ coupling is equivalent to
the category of wiring-diagram-presented interaction patterns.

OP3 is more abstract than the others. Its payoff would be a single
formal universe in which all interaction-topology choices are
expressible, which would make cross-system comparison of topology
choices easier. Its downside is that the abstraction may not be
illuminating in practice; topology choices in artificial-life
systems are often constrained by computational practicality, not by
expressive limits.

\section{The slots beyond engineered systems}

The lens was built for artificial-life systems. Most of the
philosophical heritage in Chapter~2 came from biology and the
cognitive sciences. A natural question: does the lens travel back?

Some places where the lens might be useful in non-engineered domains:

\begin{itemize}[leftmargin=1.5em,itemsep=2pt]
  \item \textbf{Cellular biology.} Cells admit a clean substrate
    (intracellular state), variation (mutation, gene expression
    noise), viability (autopoietic closure, RAF chemistry), topology
    (signalling networks), and observer (experimental measurements).
    The lens does not improve on biological models, but it might
    organise the conceptual structure differently.
  \item \textbf{Cultural evolution.} The substrate is the population
    of cultural items (memes, practices, technologies); variation is
    cultural transmission with copying error; viability is
    transmission success; topology is the social graph; the observer
    is the historian or anthropologist. The framework has a natural
    structure here, though the empirical work has barely begun.
  \item \textbf{Economic systems.} Substrate: firms, individuals,
    capital stock. Variation: innovation, hiring, investment.
    Viability: market survival. Topology: trade networks. Observer:
    GDP, employment statistics, productivity indices. The lens has
    been applied piecewise; the integrated reading does not yet
    exist.
  \item \textbf{Ecological dynamics.} Substrate: spatial-explicit
    ecology. Variation: mutation and immigration. Viability: fitness
    in habitat. Topology: dispersal kernel. Observer: species
    richness, beta diversity, food-web indices.
\end{itemize}

The framework is intentionally substrate-agnostic. Whether it
illuminates these domains, or merely re-notates them, is empirical.

\section{The bigger questions}

Beyond the technical conjectures, three larger questions linger.

\paragraph{What would a successful artificial-life system look like?}
The field has been at this for fifty years. There is no consensus
benchmark. ``Open-endedness'' is the closest thing to a target, and
even there the definitions disagree. The book has taken the cautious
position that the lens is a vocabulary, not a benchmark; success is
the ability to read and compare systems, not the production of a
single canonical run. Whether that position is good enough depends on
what one wants the field to be.

\paragraph{What is the right relationship between artificial life
and machine learning?} The two fields share variation operators
(genetic algorithms, evolution strategies) and substrates (neural
networks, LLMs) but have very different goals. Machine learning wants
solutions; artificial life wants novelty. The two have rarely been in
the same conversation. Foundation-model-mediated search (Chapter~7's
FM-embedding observers) is the first widely-adopted technique that
sits in both worlds. Whether this is a passing trend or the start of a
deeper integration is unclear.

\paragraph{Is open-endedness a real possibility?} Natural evolution
appears open-ended. Every artificial system we have built has
saturated. The conventional explanation is that our systems are too
small, too simple, or too narrowly observed. The unconventional
explanation is that open-endedness requires something we have not
identified --- a kind of variation, a kind of viability, a kind of
observer --- that the lens does not yet name. The book does not pick
between these, but a reader interested in shaping the next decade of
artificial-life research should pick.

\section{What to do next}

A few concrete starting points for a reader who wants to use the
lens.

\paragraph{Read a paper and write down its slots.} Pick a recent
artificial-life paper. Identify the substrate, the variation
operator, the viability filter, the interaction topology, and the
observer. Note which slots are reported clearly and which are
hidden. Note which conjectures the paper implicitly assumes (C1, C2,
C3 are common).

\paragraph{Run a small ablation.} Pick a system you can implement on
a laptop. Lenia + a simple novelty observer is enough. Implement it
in the lens's typed form. Vary one slot. Observe what changes.

\paragraph{Try to settle a piece of an open problem.} OP1 is the most
tractable. Pick a small system. Build its Markov chain. Search for
high-EI coarse-grainings. Even a negative result is useful: it
constrains the kinds of systems that can be causally emergent.

\paragraph{Apply the lens to a non-engineered system.} Cellular
biology, cultural evolution, or ecology. Write down the slots.
Compare with the engineered systems we have studied. Note what is
the same and what is different.

\section*{To think about}

\begin{enumerate}
  \item Pick one of C1, C2, C3 and sketch what a proof would look
    like. Sketch what a counterexample would look like. Which seems
    more likely?
  \item OP1 is the largest open problem in the framework. What
    specific candidate system would you bet on for the first
    affirmative answer? Why?
  \item The lens travels well to cellular biology, cultural
    evolution, and ecology by the book's claim. Pick a fourth
    non-engineered domain. Sketch the slot decomposition.
  \item Open-endedness is a moving target. Suppose you had a system
    that, by some operationalisation, was open-ended for a thousand
    years. What further evidence would you want before believing the
    result was real?
  \item Suppose the lens is wrong --- the five slots are not the
    right decomposition. What is the most natural alternative? What
    does it lose?
\end{enumerate}
